% §4 Stopping Rules for Outcome Prediction
The grammar of Section~\ref{sec:grammar} turns a conversation into a walk over $Q$ that emits typed observations at every turn.
A decision layer must convert this growing evidence into an action --- intervene, escalate, flag, or wait --- at a time that is early enough to matter and honest about its own error rates.
This section defines that layer: an outcome-class posterior maintained per turn, a sequential stopping rule over the posterior, and cost-asymmetric decision boundaries derived from the banking cost model.
The centerpiece is the Trajectory Stopping Theorem, which bounds the expected decision time at a prescribed confidence and subsumes the fixed-threshold triggering of learned early-failure detectors as a special case.

\subsection{From Grammar Emissions to an Outcome Posterior}
\label{subsec:posterior}
At turn $k$ the monitor observes the multi-channel tuple $o_k$ of Definition~\ref{def:observation}: the partial transcript, the turn-state sequence, the tool-event log, and the acoustic flags.
Together these channels induce the outcome-class posterior $p(\theta \mid o_{1:k})$ over the terminal set $T = \{\theta_{cont}, \theta_{aban}, \theta_{hand}, \theta_{esc}\}$.
The posterior is deliberately abstract: any calibrated estimator over grammar emissions suffices, and every guarantee below attaches to the stopping layer rather than to a particular model.
That partial trajectories carry outcome signal at all is established for offline evaluation pruning \citep{earlyeval2026} and for text-chat failure forecasting \citep{nourbakhsh2026tracer}; our contribution is a live decision layer with guarantees, not another estimator.
Two design constraints come from our published results rather than from new evidence.
First, text-only turn signals are near chance on Spanish hard cuts \citep{garcia2026asking}, so the posterior must condition on the multi-channel observation $o_{1:k}$, not on the transcript alone.
Second, the cycle depth $D(\sigma)$ of the grammar is its drift statistic, and it enters the posterior as a covariate: each regression transition shifts probability mass toward the action classes.
Deep tool episodes are precisely where banking conversations drift, and the $k$-step analysis of cascade agents makes the cost of such episodes formal \citep{garcia2026latencyfloor}.

\subsection{Sequential Stopping over the Posterior}
\label{subsec:rule}
The sequential probability ratio test (SPRT) is the canonical likelihood-ratio procedure with type-I and type-II error control \citep{wald1947sprt}, recently transposed to early-stopping of LLM reasoning paths \citep{consol2025}.
We transpose it once more, from reasoning paths to conversation turns.

\begin{definition}[Sequential outcome decision rule]
\label{def:rule}
Fix the benign class $\theta_0 = \theta_{cont}$ and, for each action class $\theta_a \in T \setminus \{\theta_0\}$, let
\begin{equation}
  \label{eq:llr}
  \ell_k^{a} \;=\; \log \frac{p(o_{1:k} \mid \theta_a)}{p(o_{1:k} \mid \theta_0)}
\end{equation}
denote the log-likelihood ratio of the observed emission stream under $\theta_a$ versus $\theta_0$.
With target false-alarm rate $\alpha_a$ and miss rate $\beta_a$ for class $\theta_a$, set the upper and lower boundaries
\begin{equation}
  \label{eq:boundaries}
  A_a = \log \frac{1}{\alpha_a}, \qquad B_a = \log \beta_a .
\end{equation}
The decision time is
\begin{equation}
  \label{eq:kdec}
  k_{\mathrm{dec}} \;=\; \min\Bigl(\;
    \min\{k : \exists\, a,\; \ell_k^{a} \geq A_a \;\text{(declare } \theta_a)\},\;\;
    \min\{k : \forall\, a,\; \ell_k^{a} \leq B_a \;\text{(declare } \theta_0)\},\;\;
    k_{\mathrm{term}}
  \Bigr),
\end{equation}
where $k_{\mathrm{term}}$ is the turn at which the grammar walk reaches a terminal state and the outcome is observed directly.
\end{definition}

The rule is one-sided per action class and jointly conservative: declaring \emph{no action} requires every class-wise ratio to fall below its lower boundary, so a conversation cannot be written off as contained while any action hypothesis retains evidence.
This asymmetry is deliberate and is priced in Section~\ref{subsec:economics}.
The boundary form \eqref{eq:boundaries} is the Ville-style parametrisation rather than the classical Wald form $\log\frac{1-\beta}{\alpha}$: because the rule monitors each class ratio past its own lower boundary (another class may still be above its $B$), the operative error event for declaring $\theta_a$ under $\theta_0$ is \emph{ever crossing} $A_a$, and Ville's inequality prices exactly that event, at the cost of slightly wider boundaries than the two-sided classical SPRT.

\subsection{The Trajectory Stopping Theorem}
\label{subsec:theorem}
We now state the guarantee.
The assumptions are explicit because each is load-bearing.
(i) \emph{Correct specification}: the emission densities used in \eqref{eq:llr} coincide with the true data-generating densities given the terminal class.
(ii) \emph{Increment separation}: emissions are conditionally independent given the terminal class, the per-turn log-ratio increments are integrable and bounded, $|Z_k^{a}| \leq c$ for a known constant $c$, and their conditional means are strictly separated: under $\theta_a$ at least $\mu_1^{a} > 0$, and under $\theta_0$ at most $\mu_0^{a} < 0$, for every action class $a$.
(iii) \emph{Termination}: the grammar walk reaches a terminal state in bounded expected time, $\mathbb{E}[k_{\mathrm{term}}] < \infty$; this is a property of the orchestration policy (Definition~\ref{def:run} makes non-terminating runs a policy bug), not of the theorem --- note that the \texttt{tool-call} self-loop makes unbounded runs representable, so the assumption must be checked against the deployed policy.

\begin{theorem}[Trajectory Stopping]
\label{thm:stopping}
Under conditions (i)--(iii), the decision rule of Definition~\ref{def:rule} satisfies, for every action class $a$:
\begin{align}
  \Pr\bigl[\text{declare } \theta_a \mid \theta_0\bigr] &\leq \alpha_a,
  &
  \Pr\bigl[\text{declare } \theta_0 \mid \theta_a\bigr] &\leq \beta_a,
  \label{eq:errors}\\
  \mathbb{E}\bigl[k_{\mathrm{dec}} \mid \theta_a\bigr]
    &\leq \frac{A_a + c}{\mu_1^{a}},
  &
  \mathbb{E}\bigl[k_{\mathrm{dec}} \mid \theta_0\bigr]
    &\leq \min\Biggl(\mathbb{E}\bigl[k_{\mathrm{term}}\bigr],\;
      \sum_{j \in A} \frac{|B_j| + c}{|\mu_0^{\,j}|}\Biggr),
  \label{eq:times}
\end{align}
where $A = T \setminus \{\theta_0\}$ is the action-class set.
The $\theta_0$-side bound aggregates over classes because declaring \emph{no action} requires
\emph{every} class-wise ratio to fall below its own lower boundary (Definition~\ref{def:rule}),
so $k_{\mathrm{dec}}$ under $\theta_0$ is gated by the slowest class-wise descent, not by any
single class.
Consequently, for any prior $\pi$ over $T$,
\begin{align}
  \label{eq:unconditional}
  \mathbb{E}[k_{\mathrm{dec}}] \;=\; \textstyle\sum_{\theta \in T} \pi_\theta\, \mathbb{E}[k_{\mathrm{dec}} \mid \theta] \;\leq\; K(\varepsilon, c),
\\[2pt]
  K(\varepsilon, c) \;:=\; \max\Biggl(\max_{a \in A} \tfrac{A_a + c}{\mu_1^{a}},\;
  \min\bigl(\mathbb{E}[k_{\mathrm{term}}], \textstyle\sum_{j \in A} \tfrac{|B_j| + c}{|\mu_0^{\,j}|}\bigr)\Biggr),
\end{align}
and by the union bound the error probability against the budgeted error events obeys
$\varepsilon := \sum_{a} (\alpha_a + \beta_a)$, so decisions are delivered in at most
$K(\varepsilon, c)$ expected turns at confidence at least $1 - \varepsilon$ against those events,
under a calibrated emission model.
A remark on the frontier: as $\beta_{esc} \to 0$ under the cost ordering of Proposition~\ref{prop:asymmetric}, $|B_{esc}|$ inflates the $\theta_0$-side sum in \eqref{eq:times} --- the price of "contained must be earned" is a slower earned verdict, an inherent SPRT trade-off the protocol measures rather than hides.
\end{theorem}

\begin{proof}[Proof sketch]
Fix an action class $a$.
Under $\theta_0$ and correct specification, the likelihood ratio of the emissions makes
$\exp(\ell_k^{a})$ a mean-one non-negative martingale: conditional independence is not needed
here, only that $\exp(\ell_k^{a})$ is the density ratio of the true conditional law, so
$\mathbb{E}_{\theta_0}[e^{Z_k^{a}} \mid \mathcal{F}_{k-1}] = 1$ exactly.
Ville's inequality then prices the ever-crossing event directly:
$\Pr_{\theta_0}[\exists k \leq k_{\mathrm{term}}: \ell_k^{a} \geq A_a] \leq e^{-A_a} = \alpha_a$,
which is the first bound of \eqref{eq:errors}; the second follows symmetrically by applying
Ville's inequality to $\exp(-\ell_k^{a})$ under $\theta_a$, whose crossing level $-B_a$ yields
$e^{B_a} = \beta_a$.
For the $\theta_a$-side time bound, use $a$'s \emph{own} upper crossing as the envelope: on the
event that $a$ is declared, $k_{\mathrm{dec}}$ is at most the first time $\ell_k^{a}$ reaches
$A_a$, and Wald's first identity (which presupposes $\mathbb{E}[k_{\mathrm{dec}}] < \infty$;
truncating at $k_{\mathrm{term}}$ and passing to the limit discharges this under assumption (iii))
gives $\mathbb{E}[\ell_k^{a} \wedge A_a] \leq A_a + c$ after $\lceil (A_a+c)/\mu_1^{a} \rceil$
turns in expectation, with Lorden-type overshoot bounds keeping the excess within $c$
\citep{lorden1970excess}.
Under $\theta_0$, declaring $\theta_0$ requires all $|A|$ ratios to descend below their $B_j$;
the expected time for the slowest descent is bounded by the sum of the per-class expected
descent times, $\sum_j (|B_j| + c)/|\mu_0^{\,j}|$, and the observed-termination mode is bounded
by $\mathbb{E}[k_{\mathrm{term}}]$ via assumption (iii), giving the second display of
\eqref{eq:times}.
The unconditional bound \eqref{eq:unconditional} takes the max of the conditional envelopes,
and the error budget is the union bound over the per-class error events
\eqref{eq:errors}.
\end{proof}

\begin{remark}[What the budget prices]
\label{rem:budget_scope}
The budget $\varepsilon = \sum_a (\alpha_a + \beta_a)$ prices two error families: declaring an
action class under the benign truth ($\theta_a \mid \theta_0$) and writing the conversation off
under an action truth ($\theta_0 \mid \theta_a$).
It does not price \emph{cross-action} misdeclarations (declaring $\theta_b$ when the truth is
$\theta_a$, $b \neq a$): $e^{\ell^b}$ is not a martingale under $\theta_a$, so Ville on
$\ell^b$ alone says nothing there; the pairwise likelihood-ratio martingale $e^{\ell^b - \ell^a}$
gives $\Pr_{\theta_a}[\text{declare } \theta_b] \leq \alpha_b / \alpha_a$ when $A_b > A_a$.
Since handoff-versus-escape confusion is economically real (Section~\ref{subsec:economics}), a
deployment that must price it adds pairwise boundaries on $\ell^a - \ell^j$ with a per-pair
budget; we leave that extension to the protocol of Section~\ref{sec:evaluation} as an operating
curve, not a theorem.
\end{remark}

Three remarks on what the theorem does and does not say.
First, the assumption we cannot discharge is per-turn conditional independence of emissions; it holds to the extent that the grammar's typed transitions, not hidden user state, drive the observation stream, and Section~\ref{sec:discussion} returns to this.
Second, the $c$ terms are the standard excess-over-boundary (overshoot) constants; in the idealised continuous treatment they vanish, and as design guidance they are one turn at most (the protocol of Section~
ef{sec:evaluation} measures them), which is why \eqref{eq:times} is displayed with them explicit rather than hidden.
Third, where the emission likelihood is unknown and the posterior is learned, likelihood ratios are unavailable in exact form.
The honest sequential alternative is an anytime-valid confidence sequence over the learned posterior, which retains error control under continuous monitoring and optional stopping \citep{howard2021time} at the price of wider early boundaries; we adopt this variant in the evaluation protocol of Section~\ref{sec:evaluation} and report both operating curves.
Wald-type SPRT remains the reference point because its boundaries are explicit and its expected stopping time is optimal among two-hypothesis tests with the same error rates \citep{waldwolfowitz1948optimality}; that optimality is a property of the classical two-sided test, not of the wider-boundary three-exit rule deployed here.

\begin{figure}[t]
  \centering
  \resizebox{\linewidth}{!}{%
    % fig_stopping.tex — stopping rule timeline over the outcome-class posterior
\begin{tikzpicture}[xscale=1.05, yscale=1.0]
% time axis: turns
\draw[-Stealth, thick] (0,0) -- (11,0) node[below right, font=\small] {turn $k$};
\foreach \k/\lbl in {1/{auth}, 3/{intent}, 5/{tool-loop}, 7/{clarify}, 9/{confirm}} {
  \draw (\k,0.08) -- (\k,-0.08) node[below, font=\scriptsize] {\lbl};
}
% posterior curves (stylized) for two classes
\draw[thick, blue!70!black] (0,1.0) .. controls (2,1.15) and (4,1.6) .. (6,2.3)
  .. controls (8,2.9) and (9.4,3.3) .. (10.6,3.45)
  node[right, font=\scriptsize, blue!70!black] {$P(\textit{contained}\mid x_{1:k})$};
\draw[thick, red!70!black] (0,1.0) .. controls (2,0.95) and (3.5,1.05) .. (5,1.5)
  .. controls (6.5,2.0) and (8.5,2.6) .. (10.6,3.1)
  node[right, font=\scriptsize, red!70!black] {$P(\textit{handoff}\mid x_{1:k})$};
% SPRT boundaries
\draw[dashed, thick, gray] (0,3.3) -- (10.6,3.3) node[above, font=\scriptsize, pos=0.85] {accept $H_1$ ($\geq 1-\delta$)};
\draw[dashed, thick, gray] (0,0.5) -- (10.6,0.5) node[below, font=\scriptsize, pos=0.85] {accept $H_0$};
% crossing point: contained class crosses upper boundary at k*
\fill[blue!70!black] (7.4,3.3) circle (2.2pt);
\draw[-Stealth, blue!70!black] (7.4,3.3) -- (7.4,0) node[below, font=\scriptsize, blue!70!black] {$k_{\mathrm{dec}}$};
% intervention window vs escalation-only threshold
\draw[decorate, decoration={brace, mirror, amplitude=5pt}, thick, olive!70!black] (7.4,4.1) -- (10.6,4.1)
  node[midway, above, font=\scriptsize, olive!70!black] {intervention window};
\node[font=\scriptsize, olive!70!black, align=left, anchor=west] at (0.1,4.55)
  {decision must land early enough to act: $t_{k_{\mathrm{dec}}} + \varepsilon$ before the escalation-only threshold, and the intervention itself must clear the cessation floor};
\end{tikzpicture}%
  }
  \caption{The stopping layer over a live conversation.
           Class-wise log-likelihood ratios $\ell_k^{a}$ accumulate with their boundary pair
           $(B_a, A_a)$; the first upper crossing declares an action class and fires an
           intervention whose decision latency $\lambda_{\mathrm{decide}}$ and cessation
           latency $\lambda_{\mathrm{cease}}$ must both fit inside the remaining window
           before the escalation-only threshold.}
  \label{fig:stopping}
\end{figure}

\subsection{Floor Composition: Early Enough to Act}
\label{subsec:floors}
A turn index is not an intervention guarantee, so we convert $k_{\mathrm{dec}}$ to wall-clock using the published floor arithmetic.
The per-turn decision latency $\lambda_{\mathrm{decide}}$ of the monitor must fit inside the inter-turn budget that the onset floor leaves unfilled \citep{garcia2026latencyfloor}, and the measured decision latencies of the Spanish interim-transcript pipeline (p99 178--268~ms, cloud) show this budget exists in production \citep{garcia2026asking}.
Executing the intervention is a second bound: taking the floor from an agent mid-utterance is a barge-in, and it must clear the cessation floor's detect--decide--cease--resume decomposition \citep{garcia2026cessationfloor}.
The composition is therefore
\begin{equation}
  \label{eq:compose}
  t_{k_{\mathrm{dec}}} + \lambda_{\mathrm{decide}} + \lambda_{\mathrm{cease}}
  \;<\;
  t_{\mathrm{escalation\text{-}only}},
\end{equation}
and only conversations satisfying \eqref{eq:compose} are ones where early prediction converts into an actionable intervention rather than a post-hoc annotation.
The cycle-length bound of Section~\ref{sec:grammar} completes the picture from the other side: a run segment exceeding $k_{\mathrm{cyc}}$ \texttt{tool-call} iterations violates the response deadline, so the stopping rule rarely has to wait past the horizon the grammar already justifies.

\subsection{False-Early-Exit Economics}
\label{subsec:economics}
Stopping early is only justified if the cost model says so.
Let $c(\hat{\theta}, \theta)$ denote the banking cost of declaring class $\hat{\theta}$ when the truth is $\theta$, over declared versus true terminal classes; its operational content is specified with the KPI wiring of Section~\ref{sec:evaluation}.
False containment --- declaring \textit{contained} while the truth is \textit{handoff} or \textit{compliance-escape} --- is the expensive error: it converts a recoverable failure into an unattended one, and a compliance escape has unbounded cost.
False escalation --- declaring \textit{handoff} while the truth is \textit{contained} --- wastes teller minutes, the trade that queue-level control already prices from the outside of the conversation \citep{queuecontrol2026}.
Abandonment adds a patience-driven cost that queueing models estimate from uncertain data \citep{silentabandonment2024}; the grammar lets us price it from the conversation's own trajectory.

\begin{proposition}[Cost-asymmetric boundaries]
\label{prop:asymmetric}
Let the cost model satisfy $c(\theta_0, \theta_a) > c(\theta_a, \theta_0)$ for every action class $a$, with $c(\theta_0, \theta_{esc})$ unbounded.
Then the Bayes-optimal single-turn threshold for declaring $\theta_{esc}$ tends to zero, so the compliance boundary fires at the first credible signal: $\alpha_{esc}$ is elevated, $A_{esc} = \log(1/\alpha_{esc})$ is the smallest upper boundary, and $\beta_{esc}$ tends to zero.
Between the remaining action classes, larger relative miss cost $c(\theta_0, \theta_a)/c(\theta_a, \theta_0)$ yields a smaller threshold, hence a smaller $A_a$ and a larger $\alpha_a$.
\end{proposition}

\begin{proof}[Proof sketch]
For a single-turn Bayes decision with posterior $p(\theta_a \mid o_{1:k})$, the optimal rule declares $\theta_a$ exactly when the posterior exceeds the cost ratio $\frac{c(\theta_a,\theta_0)}{c(\theta_a,\theta_0) + c(\theta_0,\theta_a)}$.
An unbounded $c(\theta_0, \theta_{esc})$ drives that ratio to zero, so the compliance escape fires at any positive posterior mass; in the Ville parametrisation \eqref{eq:boundaries} this is an elevated false-alarm rate $\alpha_{esc}$ and hence the smallest upper boundary $A_{esc}$, while the miss rate $\beta_{esc}$ --- the probability of declaring $\theta_0$ when the truth is $\theta_{esc}$ --- is driven to zero by the same eagerness.
Monotonicity in the relative miss cost follows from differentiating the threshold ratio.
The sequential rule inherits the ordering because its boundaries are the sequential realization of the same posterior-threshold test at every turn.
\end{proof}

The practical reading is that the four terminal classes do not share one threshold.
A banking monitor should fire its compliance escape at the first credible signal, accept a higher false-escape rate as the price of unbounded miss cost, and treat contained as a verdict that must be earned against every action hypothesis simultaneously.

\subsection{TRACER as the Fixed-Sample Instance}
\label{subsec:tracer}
The nearest published system to this decision layer is TRACER, which scores every visible dialogue prefix for eventual failure and fires a recovery trigger when the score exceeds a development-tuned threshold $\tau = 0.7788$, frozen under a false-positive-rate constraint of $0.15$ \citep{nourbakhsh2026tracer}.

\begin{corollary}[Fixed-sample instance]
\label{cor:tracer}
TRACER's thresholded trigger is the fixed-sample instance of Definition~\ref{def:rule}: a single look at each prefix with a learned posterior, a one-sided boundary in place of the pair \eqref{eq:boundaries}, and no optional stopping.
The Trajectory Stopping Theorem adds the three properties that a single-look trigger cannot express: an expected decision-time bound \eqref{eq:times}, error control that survives continuous monitoring (via the confidence-sequence variant of \citealp{howard2021time}), and cost-derived asymmetry across the four banking outcome classes instead of one binary failure label.
\end{corollary}

This is a generalization claim, not a criticism: TRACER's detected signals appear from the first quarter of the dialogue and its operating point is honest about its false-positive budget \citep{nourbakhsh2026tracer}.
What the sequential layer contributes is the freedom to stop looking, the price of that freedom in turns, and thresholds that differ per outcome class because the costs differ.
The evaluation protocol of Section~\ref{sec:evaluation} therefore uses TRACER's early-detection metrics with credit, extended with the cost-weighted decision-latency accounting that floor composition demands.
